\documentclass[../Main.tex]{subfiles}

\begin{document}
\section{Game Theory}
\subsection{Games and Simple Strategies}
\begin{definition}{Game}
    A \underline{game} is a setup where two players can choose from a set number of actions which affects their scores.
\end{definition}
\begin{definition}{Payoff matrix}
    For each player, the \underline{payoff matrix} is an $m \times n$ matrix whose entries $A_{ij}$ are the score (or payoff) attained by player $1$ when they play action $i$ and player 2 plays action $j$.
\end{definition}
\begin{definition}{Zero-sum game}
    A game is \underline{zero-sum} if after each round of the game the total payoff is zero. In this case the two player have payoff matrices $A$ and $-A$.
\end{definition}
\begin{definition}{Conservative approach}
    The \underline{conservative approach} for a player is to choose the action that satisfies:
    \begin{equation*}
        \max_{i \in \{1, \cdots, m\}} \min_{j \in \{1, \cdots, n\}} a_{ij}
    \end{equation*}
    which may not exist.
\end{definition}
\begin{definition}{Strategy}
    A \underline{strategy} is the probability distribution of actions that a player picks. A \underline{pure strategy} is the strategy that only ever chooses one action. Otherwise the strategy is a \underline{mixed strategy}.
\end{definition}
\subsection{Linear Programming for Strategies}
Player 1 must solve:
\begin{align*}
    \text{Maximise } &v \\
    \text{Subject to } & A^T \vec{p} \geq v\vec{e} \\
    & \vec{e} \cdot \vec{p} = 1 \\
    & \vec{p} \geq \vec{0}
\end{align*}
where $\vec{e}$ is the vector of all $1$ entries.

Player 2 solves the dual problem:
\begin{align*}
    \text{Minimise } &w \\
    \text{Subject to } & A \vec{q} \leq w\vec{e} \\
    & \vec{e} \cdot \vec{q} = 1 \\
    & \vec{q} \geq \vec{0}
\end{align*}
\begin{definition}{Nash equilibrium}
    The \underline{Nash equilibrium} is the pair $(\vec{p}, \vec{q})$ that solves the primal and dual problems.
\end{definition}
\begin{theorem}
    A strategy $\vec{p}$ is optimal for player 1 if and only if there exist $\vec{q}$ and $v$ such that the following hold.
    \begin{enumerate}
        \item Primal feasibility: $A^T\vec{p} \geq v\vec{e}, \vec{e} \cdot \vec{p} = 1, \vec{p} \geq \vec{0}$;
        \item dual feasibility: $A\vec{q} \leq v\vec{e}, \vec{e} \cdot \vec{q} = 1, \vec{q} \geq \vec0$;
        \item complementary slackness: $v = \vec{p}^TA\vec{q}$.
    \end{enumerate}
    \label{thmGameOptimality}
\end{theorem}
\section{Transport Problems}
\subsection{Network Flow Problems}
\begin{definition}{Directed graph}
    A \underline{directed graph} is a set $G = (V, E)$ where $V$ is a set of \underline{vertices}, and $E \subseteq V \times V$ is a set of \underline{edges}. The ordering of the vertices in the edge is important, since some edges have only one direction of travel.
\end{definition}
Consider a graph with $n$ vertices. Let $b_i$ represent the quantity of a good produced (if positive) or consumed (if negative) at vertex $i$. We assume that the overall amount of goods stays the same:
\begin{equation*}
    \sum_{i=1}^{n} b_i = 0.
\end{equation*}
Let $x_{ij}$ be the amount of a good flowing from $i$ to $j$ and let $C_{ij}$ be the cost per unit of a good flowing from $i$ to $j$. If there is no edge $(i, j)$ then we notationally set $C_{ij}$ to be undefined or infinite. 

We also have matrices $\underline{M}, \overline{M}$ such that:
\begin{equation*}
    \underline{M}_{ij} \leq x_{ij} \leq \overline{M}_{ij}
\end{equation*}
\begin{definition}{Network flow problem}
    A \underline{network flow problem} has the form:
    \begin{align*}
        \text{Minimise } & \sum_{(i, j) \in E} C_{ij} x_{ij} \\
        \text{Subject to } & b_i + \sum_{j : (j, i) \in E} x_{ji} = \sum_{j : (i, j) \in E} x_{ij} ~\forall i \in V \\
        & \underline{M}_{ij} \leq x_{ij} \leq \overline{M}_{ij}~\forall (i, j) \in E
    \end{align*}
\end{definition}
\subsection{Transport Problem}
\begin{definition}{Transport problem}
    A \underline{transport problem} is a network flow problem where each vertex is either a producer or consumer, and goods flow only from producers to consumers.
\end{definition}
Consider $n$ producers and $m$ consumers. Let producer $i$ supply $s_i$ of a good and consumer $j$ demand $d_j$ of it.
\begin{definition}{Transport problem}
    A \underline{transport problem} has the form:
    \begin{align*}
        \text{Minimise } & \sum_{(i, j) \in E} C_{ij} x_{ij} \\
        \text{Subject to } & \sum_{j=1}^{m} x_{ij} = s_i~\forall 1 \leq i \leq n \\
        & \sum_{i=1}^{n} x_{ij} = d_j~\forall 1 \leq j \leq m
    \end{align*}
\end{definition}
\begin{theorem}
    Any network flow problem with finite capacities or non-negative costs has an equivalent transport problem.
    \label{thmTransportNetFlowEquiv}
\end{theorem}
\begin{proof}
    First observe that $\underline{M}$ can be set to $0$ by re-labelling $x_{ij}$ to $\tilde{x}_{ij} - \underline{M}_{ij}$.

    Define also:
    \begin{equation*}
        \tilde{b}_i = \sum_{(i, j)\in E} \underline{M}_{ij}.
    \end{equation*}
    Replace any infinite-capacity edge with the finite (but unattainable) capacity $\overline{M}_{ij} = \sum_i \abs{b_i}$.

    For every edge $(i, j) \in E$ construct a supplier and for every vertex construct a consumer. Let the supply from supplier $(i, j)$ be the capacity $\overline{M}_{ij}$. Let the cost to supply $i$ be zero, and to supply $j$ be $c_{ij}$, with no other connections.

    Set the amount of goods consumed by consumer $i$ to:
    \begin{equation*}
        \sum_{k : (i, k) \in E} \overline{M}_{ik} - b_i
    \end{equation*}
    and by consumer $j$ to:
    \begin{equation*}
        \sum_{k : (k, j) \in E} \overline{M}_{jk} - b_j.
    \end{equation*}
    Now send a flow $x_{ij}$ from producer $(i, j)$ to consumer $j$. Then the flow from this producer to consumer $i$ must be $\overline{M}_{ij} - x_{ij}$.

    Note that the cost function is therefore the same, and constructs the constraints and show these are the same too. Construct the inflow~$=$~demand equality, and simplify it to show it is equivalent to that of the original problem.
\end{proof}
\subsection{Transport Algorithm}
\begin{theorem}[Transport Optimality Theorem]
    For a feasible flow $x$, if there exists $\vec{\lambda} \in \R^n$ and $\vec{\mu} \in \R^m$ such that, for all $i$ and $j$,
    \begin{align*}
        c_{ij} &\geq \lambda_i + \mu_j \\
        0 &= (c_{ij} - (\lambda_i + \mu_j))x_{ij}
    \end{align*}
    then $x$ is optimal.
    \label{thmTransportOptimal}
\end{theorem}
\begin{proof}
    The proof is a specialisation of that of \thmref{thmLinearProgOptimality}.
\end{proof}
We perform the transport algorithm on a Tableau which keeps track of $\vec{\lambda}, \vec{\mu}$ and $x$.

\begin{tabular}{c|c|c|c|c|c}
     & $\mu_1$ & $\mu_2$ & & $\mu_m$ & \\
    \hline
    $\lambda_1$ & \transportcell{$\lambda_1 + \mu_1$}{$x_{11}$}{$c_{11}$} & \transportcell{$\lambda_1 + \mu_2$}{$x_{12}$}{$c_{12}$} & $\cdots$ & \transportcell{$\lambda_1 + \mu_m$}{$x_{1m}$}{$c_{1m}$} & $s_1$ \\
    \hline
    $\lambda_2$ & \transportcell{$\lambda_2 + \mu_1$}{$x_{21}$}{$c_{21}$} & \transportcell{$\lambda_2 + \mu_2$}{$x_{22}$}{$c_{22}$} & $\cdots$ & \transportcell{$\lambda_2 + \mu_m$}{$x_{2m}$}{$c_{2m}$} & $s_2$ \\
    \hline
    $\vdots$ & $\vdots$ & $\vdots$ & $\ddots$ & $\vdots$ & $\vdots$ \\
    \hline
    $\lambda_n$ & \transportcell{$\lambda_n + \mu_1$}{$x_{n1}$}{$c_{n1}$} & \transportcell{$\lambda_n + \mu_2$}{$x_{n2}$}{$c_{n2}$} & $\cdots$ & \transportcell{$\lambda_n + \mu_m$}{$x_{nm}$}{$c_{nm}$} & $s_n$ \\
    \hline
     & $d_1$ & $d_2$ & $\cdots$ & $d_m$

\end{tabular}

We must initialise the algorithm with a basic solution. This can be constructed using the north-west corner rule:
\begin{enumerate}
    \item start with the first supplier and first consumer;
    \item push as much of the good along this edge as possible (either running out of supply or demand);
    \item if supply was exhausted, move to the next supplier;
    \item if demand was exhausted, move to the next consumer;
    \item repeat from step 2.
\end{enumerate}
This is guaranteed to generate a feasible solution, except in the case of a degenerate problem.

To perform one step of the algorithm, initially set $\lambda_1$ to zero and use complementary slackness to determine the rest of $\vec{\lambda}$ and $\vec{\mu}$. If then dual feasibility holds, stop.

If dual feasibility does not hold for some $(i, j)$, push an amount $\theta$ along this new edge. Continue around adding and subtracting $\theta$ to ensure feasibility, until reaching a loop. Choose the largest $\theta$ such that no edge has negative flow (thereby deleting an edge).

This is the next feasible solution to check and then perform another step of the algorithm if not yet optimal.
\section{Max-Flow Min-Cut}
\begin{definition}{Max-flow min-cut problem}
    A \underline{max-flow min-cut problem} is a network flow problem with a single source at vertex $1$ and a single sink at vertex $n$. It has the form:
    \begin{align*}
        \text{Maximise } &\delta \\
        \text{Subject to } & \sum_{j : (1, j) \in E} x_{1j} - \sum_{j : (j, 1) \in E} x_{j1} = \delta \\
        & \sum_{j : (n, j) \in E} x_{nj} - \sum_{j : (j, n) \in E} x_{jn} = -\delta \\
        & \sum_{j : (i, j) \in E}x_{ij} - \sum_{j : (j, i) \in E} x_{ji} = 0,~~i \neq 1, n \\
        &0 \leq x_{ij} \leq c_{ij}
    \end{align*}
\end{definition}
\begin{definition}{Graph cut}
    A \underline{cut} of a graph $G = (V, E)$ is a partition of its vertices into $S$ and $V \backslash S$ such that $1 \in S, n \notin S$.
\end{definition}
\begin{definition}{Cut capacity}
    The \underline{capacity} of a graph cut is:
    \begin{equation*}
        C(S) = \sum_{\substack{(i, j) \in E\\i \in S\\j \notin S}} c_{ij}
    \end{equation*}
\end{definition}
\begin{proposition}
    Let $x$ be a feasible flow. Then for any cut $S$, $C(S) \geq \delta$.
    \label{propMaxFlowAtLeastMinCut}
\end{proposition}
\begin{proof}
    First define the flow from a set $X$ to $Y$, not necessarily disjoint:
    \begin{equation*}
        f_x(X, Y) = \sum_{\substack{(i, j) \in E\\i \in X\\j \in Y}} c_{ij}
    \end{equation*}
    Sum up all the flow constraints over the vertices to get:
    \begin{equation*}
        \delta = f_x(S, V) - f_x(V, S)
    \end{equation*}
    and by splitting $V$ into its partition, show that also:
    \begin{equation*}
        \delta = f_x(S, V\backslash S) - f_x(V\backslash S, S) \leq C(S) + 0
    \end{equation*}
\end{proof}
\begin{theorem}[Max-flow min-cut theorem]
    The value of the maximum flow is the value of the minimum cut capacity.
    \label{thmMaxFlowMinCut}
\end{theorem}
\begin{proof}
    Define an augmenting path as a sequence of connected vertices along which the flow can be increased. Argue that if an augmenting path from $1$ to $n$ exists then we do not have optimality.

    Define $S$ as the set of vertices that can be reached from $1$ by an augmenting path. In the optimal case $n \notin S$ and so $S$ defines a cut. In the proof of \thmref{propMaxFlowAtLeastMinCut} we see:
    \begin{equation*}
        \delta = f_x(S, V\backslash S) - f_x(V\backslash S, S)
    \end{equation*}
    where the first term is $C(S)$, and we can show that the second term is $0$ by contradiction.
\end{proof}
The Ford-Fulkerson Algorithm uses this idea:
\begin{enumerate}
    \item start with a feasible flow (such as zero everywhere);
    \item find an augmenting path from $1$ to $n$;
    \item push the max flow along it;
    \item repeat until there are no more augmenting paths.
\end{enumerate}
This will find a solution for integer capacities, but may not find a solution otherwise.
\end{document}